%%%%%%%%%%%%%%%%%%%%%%%%%%%%%%%%%%%%%%%%%%%%%%%%%%%%%%%%%%%%%%%%%%
%% Toto je inkludovaný LaTeXový soubor `ca_lopkg.tex'.          %%
%% ------------------------------------------------------------ %%
%% Verze 2022-06-03                                             %%
%% Vyvíjí & udržuje <stanislav.hledik@physics.slu.cz>           %%
%%%%%%%%%%%%%%%%%%%%%%%%%%%%%%%%%%%%%%%%%%%%%%%%%%%%%%%%%%%%%%%%%%

\chapter{Načítané balíčky}\label{packages}%%%%%%%%%%%%%%%%%%%%%%%%%%%%%%%%%

Následující balíčky jsou načítány v~uvedeném pořadí dokumentní třídou
\pgm{\ipotname .cls} pomocí příkazu \pgm{\textbackslash RequirePackage}, takže
není třeba je v~preambuli pomocí příkazu \pgm{\textbackslash usepackage} znovu
načítat.\footnote{Balíčky \pkg{hyperref} resp. \pkg{url} jsou načítány
  v~argumentu makra \protect\pgm{\textbackslash AtEndPreamble}, což zajistí
  načtení \pkg{hyperref} jako posledního i~v~případě, že v~preambuli budou
  načítány další uživatelem požadované balíčky.}

\begin{description}
  \pkgitem{etoolbox}\label{pkgetoolbox} je sada programovacích nástrojů
  primárně určených pro autory \LaTeX ových tříd a
  balíčků~\cite{Wri-Leh:LIVE:CTAN:etoolbox}. Díky makru \pgm{\textbackslash
    AtEndPreamble} definovanému tímto balíčkem je možné odložit načtení
  určitého kódu, včetně balíčku \pkg{hyperref}, až na samý konec preambule,
  nikoli však do argumentu \pgm{\textbackslash AtBeginDocument}, kde by nové
  verze \pkg{hyperref} nefungovaly správně.
  
\pkgitem{natbib}\label{pkgnatbib} podporuje flexibilitu vytváření referencí
  a citací pomocí pomocného programu \BibTeX{}~\cite{Dal:LIVE:CTAN:natbib,%
    Dal:LIVE::Natbib,%
    Kop-Dal:2004:LaTeXPodrPruv:,%
    Goo-Mit-Sam:1994:Companion:}.
% Pro formátování databázových záznamů do \LaTeX u\INDEX{publikační
% systém!\LaTeX{}} používá \RINDEX{bibliografický styl} \pgm{\ipotname
% .bst}, jenž je součástí balíčku \pkg{\ipotname}.
%  (viz Dodatek~\ref{files}).
%   Pro jeho správnou funkci je třeba mít ve vyhledávacích cestách soubor
%   \pgm{babelbst.tex} a lokalizační \pgm{*.bld} soubory pro jednotlivé jazyky
%   (vše je součástí \pkg{\ipotname} s~podporou jazyků: čeština, angličtina
%   britská a americká, slovenština, francouzština, němčina).
% %  (viz též sekci~\ref{bibliodb}).

\pkgitem{amsmath, amssymb}\label{pkgams} výrazně rozšiřují možnosti \LaTeX
  u\INDEX{publikační systém!\LaTeX{}} v~oblasti sazby matematiky jak co do
  definic příkazů a prostředí, tak výběru
  symbolů~\cite{LtX:LIVE:CTAN:AMSmath,%
    AMS:LIVE:CTAN:AMSfonts,%
    Kop-Dal:2004:LaTeXPodrPruv:,%
    Goo-Mit-Sam:1994:Companion:}.\footnote{Jelikož \pkg{amssymb} interně
    načítá \pkg{amsfonts} (jak se lze přesvědčit nahlédnutím do logu
    překladu), není již třeba \pkg{amsfonts} explicitně načítat. Zamýšlíte-li
    užívat matematický styl definice--věta--důkaz, zvažte načtení balíčku
    \pgm{amsthm} v~preambuli.}

\pkgitem{graphicx}\label{pkggraphicx} poskytuje, mimo jiné, unifikované
  rozhraní pro vkládání obrázků obsažených v~externích
  souborech~\cite{Car:LIVE:CTAN:Graphics,%
    Kop-Dal:2004:LaTeXPodrPruv:,%
    Goo-Rah-Mit:1994:GraComp:}.

\pkgitem{hyperref, url}\label{pkghyperref} slouží k~transformaci externích
  i~interních křížových odkazů na
  hypertext~\cite{LtX:LIVE:CTAN:hyperref}. Balíček~\pkg{url} nutný pro řádkový
  zlom URL v~místě teček, lomítek
  atd~\cite{Ars:LIVE:CTAN:url,Goo-Mit-Sam:1994:Companion:} se implicitně načítá
  přes balíček \pkg{hyperref}, jenž je sám načítán pouze při zapnuté volbě
  \opt{htext} (viz stranu~\pageref{opthtext}); v~opačném případě je tedy nutno
  jej načíst samostatně~-- vše je ošetřeno v~\pgm{\ipotname .cls}.
\end{description}

\noindent
Máte-li v~úmyslu načítat v~preambuli nějaké balíčky, vyhněte se opětovnému
načítání výše vyjmenovaných nebo v~sekci~\ref{dependencies} níže uvedených;
mohlo by to vést k~těžko odhalitelným chybám. Obzvláště citlivý je na to
balíček \pkg{hyperref}~\cite{LtX:LIVE:CTAN:hyperref}, u~kterého je automaticky
zajištěno načtení jako posledního ze všech. Například potřebujeme-li pro
matematicky zaměřenou závěrečnou práci balíček
\pkg{amsthm}~\cite{AMS:LIVE:CTAN:AMSthm}, stačí v~preambuli načíst pouze jej,
jelikož na něj úzce vázané balíčky \pkg{amsmath} a \pkg{amssymb}
resp. \pkg{amsfonts}~\cite{LtX:LIVE:CTAN:AMSmath,%
  AMS:LIVE:CTAN:AMSfonts} jsou již načteny.

\section{Závislosti}\label{dependencies}%%%%%%%%%%%%%%%%%%%%%%%%%%%%%%%%%%%

Každý z~uvedených balíčků závisí na množství dalších balíčků, kromě toho
v~preambuli jsou zaváděny další balíčky a ty také mohou záviset na
dalších. V~následujících seznamech jsou uvedeny pouze další balíčky načítané
těmi, jenž jsou zaváděny dokumentní třídou \pgm{\ipotname .cls}\footnote{Platí
  pro aktuální verzi \pkg{\ipotname} a distribuci \distroMiKTeX{} verze~22.3
  pro Windows~10.}  (jejich seznam je výše) a těmi, jenž jsou zaváděny
v~preambuli hlavního zdrojového souboru \pgm{\thesnamecze .tex} tohoto
manuálu\footnote{Ty se mohou měnit v~závislosti na tom, jaké balíčky načte
  v~preambuli uživatel.}:

\begin{description*}
\raggedright
%\sloppy
\item[Bez volby \opt{htext} (viz s.~\pageref{opthtext}):]
%natbib,
%amsmath,
\pkg{amstext},
\pkg{amsgen},
\pkg{amsbsy},
\pkg{amsopn},
%amssymb,
\pkg{amsfonts},
%graphicx,
\pkg{keyval},
\pkg{graphics},
\pkg{trig},
%lmodern,
%fontenc,
%inputenc,
%babel,
%fancyvrb,
%url,
\pkg{epstopdf-base} (pouze při kompilaci formátem \pkg{pdflatex})
\item[S~volbou \opt{htext}:]
%etoolbox,  
%natbib,
%amsmath,
\pkg{amstext},
\pkg{amsgen},
\pkg{amsbsy},
\pkg{amsopn},
%amssymb,
\pkg{amsfonts},
%graphicx,
\pkg{keyval},
\pkg{graphics},
\pkg{trig},
%lmodern,
%fontenc,
%inputenc,
%babel,
%fancyvrb,
%hyperref,
\pkg{ltxcmds},
\pkg{iftex},
\pkg{pdftexcmds},
\pkg{kvdefinekeys},
\pkg{hycolor},
\pkg{letltxmacro},
\pkg{auxhook},
\pkg{nameref},
\pkg{refcount},
\pkg{gettitlestring},
\pkg{kvoptions},
\pkg{intcalc},
\pkg{etexcmds},
%url,
\pkg{bitset},
\pkg{bigintcalc},
\pkg{atbegshi-ltx},
\pkg{atveryend-ltx},
\pkg{rerunfilecheck},
\pkg{uniquecounter},
\pkg{epstopdf-base} (pouze při kompilaci formátem \pkg{pdflatex}),
\pkg{color}
\end{description*}

\endinput

%%
%% End of file `ca_lopkg.tex'.
